\section{Stationary ARMA Processes}

\subsection{ARMA Processes}

\begin{definition}
    [White Noise] The process $(Z_t)$ is said to be white noise with mean $0$ and
variance $\sigma^2$ if and only if $(Z_t)$ has zero mean and covariance function
\[
\gamma(h) = \begin{cases}
    0,\quad &h \neq 0 \\ \sigma^2, & h = 0 
\end{cases}
\]
\end{definition}

\begin{definition}
    [The ARMA$(p, q)$ Process] The process  $(X_t)_{t\in\mathbb{Z}}$ is said to be an ARMA -$(p,q)$ process if $(X_t)$ is stationary and if for every $t$,
    \[
    X_t - \phi_1 X_{t-1} - \cdots \phi_p X_{t-p} = Z_t + \theta_1 Z_{t-1} + \cdots + \theta_q Z_{t-q}
    \]
where $(Z_t) \sim \text{WN}(0,\sigma^2)$. We say that $(X_t)$ is an ARMA$(p,q)$ process with mean $\mu$ if $(X_t - \mu)$ is an ARMA$(p,q)$ process.\par
The equation can be written as 
\[
\phi(B)X_t = \theta(B)Z_t
\]
where $\phi,\theta$ are $p,q$-degree polynomials and $B$ is the push-backward operator.
\end{definition}

\begin{definition}
    [The MA$(q)$ process and the AR$(p)$ process]
    The MA$(q)$ process is an ARMA$(0,q)$ process, that is $\phi \equiv 1$, i.e.
    \[
    X_t = \theta(B)Z_t.
    \]\par
    The AR$(p)$ process is an ARMA$(p,0)$ process, that is $\theta \equiv 1$, i.e.
    \[
    \phi(B)X_t = Z_t
    \]
\end{definition}

\begin{proposition}
    The mean value and covariance function for an MA$(q)$ process $X_t$ with MA polynomial $\theta$ is
    \[
    \mu_X(t) = 0,\quad \gamma_X(h) = \sigma^2\left(\sum\limits_{k=0}^{q-h}\theta_k\theta_{k+h}\right).
    \]
\end{proposition}
\begin{proof}
    We know
    \[
    \mathbb{E} X_t = \mathbb{E}\theta(B)Z_t = 0
    \]
    and
    \[
    \gamma_X(h) = \mathbb{E}\left(\mathbb{E}\theta(Z)Z_t\right)\left(\theta(Z)Z_{t+h}\right) = \sigma^2\left(\sum\limits_{k=0}^{q-h}\theta_k\theta_{k+h}\right)
    \]
    where $\theta_0:=1$.
\end{proof}

\begin{definition}[Causality]
An ARMA$(p,q)$ process defined by the equations
$\phi(B)X_t=\theta(B)Z_t$ is said to be \emph{causal}
(or more specifically to be a \emph{causal function of $\{Z_t\}$})
if there exists a sequence of constants $\{\psi_j\}_{j\geq0}$ such that
\[
\sum_{j=0}^{\infty}|\psi_j|<\infty
\]
and
\[
X_t=\sum_{j=0}^{\infty}\psi_j Z_{t-j},
\qquad t=0,\pm1,\ldots.
\]
\end{definition}

\begin{proposition}
If $\{X_t\}_{t\in\mathbb{Z}}$ is any sequence of random variables such that
$\sup_t\mathbb{E}|X_t|<\infty$, and if
$\sum_{j=-\infty}^{\infty}|\psi_j|<\infty$, then the series
\[
\psi(B)X_t
=\sum_{j=-\infty}^{\infty}\psi_j B^jX_t
=\sum_{j=-\infty}^{\infty}\psi_j X_{t-j}
\]
converges absolutely with probability one. If in addition
$\sup_t\mathbb{E}|X_t|^2<\infty$, then the series converges in mean square
to the same limit.
\end{proposition}
\begin{proof}
    To see the absolute convergence for a.s., we may notice that
    \[
    \mathbb{E} |\phi|(B)|X_t| \leq \sum\limits_{j=-\infty}^{\infty}|\phi_j|\sup\mathbb{E}|X_t| < \infty
    \]
    and hence $|\phi|(B)|X_t|$ has to be finite in probability 1.\par
    Then we would like to show that the series converge to $\phi(B)X_t$ in $L^2$ as well, that is we may prove that by noticing
    \[
    \mathbb{E}\left|\sum\limits_{m < |j| \leq n} \phi_j X_{t-j}\right|^2 \leq \left(\sum\limits_{|j|\geq m} |\phi_j|\right)^2\sup_{t}\mathbb{E}|X_t|^2 \to 0
    \]
    and hence the series converge in $L^2$ to some limit $L$, with
    \[
    \mathbb{E}|L-\sum\limits_{j = -\infty}^{\infty} \phi_jX_{t-j}|^2 \leq \liminf\mathbb{E}|L - \sum\limits_{j=-n}^n \phi_j X_{t-j}|^2 = 0
    \]
    by Fatou's lemma and we are done.
\end{proof}

\begin{proposition}
If $\{X_t\}$ is a stationary process with autocovariance function
$\gamma(\cdot)$ and if $\sum_{j=-\infty}^{\infty}|\psi_j|<\infty$,
then for each $t\in\mathbb{Z}$ the series
\[
\psi(B)X_t=\sum_{j=-\infty}^{\infty}\psi_jX_{t-j}
\]
converges absolutely with probability one and in mean square to the same
limit. If
\[
Y_t=\psi(B)X_t,
\]
then the process $\{Y_t\}$ is stationary with autocovariance function
\[
\gamma_Y(h)=\sum_{j,k=-\infty}^{\infty}
\psi_j\psi_k\gamma(h-j+k).
\]
\end{proposition}

\begin{theorem}
Let $\{X_t\}$ be an ARMA$(p,q)$ process for which the polynomials
$\phi(\cdot)$ and $\theta(\cdot)$ have no common zeroes. Then $\{X_t\}$
is causal if and only if $\phi(z)\neq0$ for all $z\in\mathbb{C}$ such that
$|z|\leq1$. The coefficients $\{\psi_j\}$ in the causal representation
$X_t=\sum_{j=0}^{\infty}\psi_jZ_{t-j}$ are determined by the relation
\[
\psi(z)=\sum_{j=0}^{\infty}\psi_jz^j
=\frac{\theta(z)}{\phi(z)},
\qquad |z|\leq 1.
\]
\end{theorem}
\begin{proof}
Suppose first that $\phi(z)\neq0$ for every $|z|\leq1$. Since the
zero set of the polynomial $\phi$ is finite, there is some $R>1$ such
that $\phi$ has no zero in $|z|<R$. Hence the rational function
\[
\psi(z):=\frac{\theta(z)}{\phi(z)}
\]
is analytic in $|z|<R$ and has a power-series expansion
\[
\psi(z)=\sum_{j=0}^{\infty}\psi_jz^j.
\]
Choose $r$ with $1<r<R$. Cauchy's estimate gives
\[
|\psi_j|\leq M_r r^{-j},
\qquad
M_r:=\max_{|z|=r}|\psi(z)|,
\]
so $\sum_{j=0}^{\infty}|\psi_j|<\infty$. Consequently
\[
Y_t:=\psi(B)Z_t=\sum_{j=0}^{\infty}\psi_jZ_{t-j}
\]
is well defined, and the preceding propositions show that $\{Y_t\}$ is
stationary. Since $\phi(z)\psi(z)=\theta(z)$, absolute summability permits
composition of the filters and gives
\[
\phi(B)Y_t=\theta(B)Z_t=\phi(B)X_t.
\]
Put $U_t=X_t-Y_t$. Then $\{U_t\}$ is stationary and
$\phi(B)U_t=0$. The function $1/\phi(z)$ is also analytic in $|z|<R$
and has absolutely summable Taylor coefficients. Applying the corresponding
filter to the last equality yields
\[
U_t=\frac{1}{\phi}(B)\phi(B)U_t=0.
\]
Thus $X_t=Y_t=\sum_{j=0}^{\infty}\psi_jZ_{t-j}$, so $\{X_t\}$ is
causal.

Conversely, suppose that $\{X_t\}$ is causal, so that
\[
X_t=\psi(B)Z_t=\sum_{j=0}^{\infty}\psi_jZ_{t-j},
\qquad
\sum_{j=0}^{\infty}|\psi_j|<\infty.
\]
Substitution into the ARMA equation gives
\[
\bigl(\phi(B)\psi(B)-\theta(B)\bigr)Z_t=0.
\]
Write
\[
\phi(z)\psi(z)-\theta(z)=\sum_{j=0}^{\infty}a_jz^j.
\]
The sequence $\{a_j\}$ is absolutely summable. Since $\{Z_t\}$ is
white noise with variance $\sigma^2>0$,
\[
0=\operatorname{Var}\left(\sum_{j=0}^{\infty}a_jZ_{t-j}\right)
=\sigma^2\sum_{j=0}^{\infty}|a_j|^2.
\]
Therefore $a_j=0$ for every $j$, and hence
\[
\phi(z)\psi(z)=\theta(z),
\qquad |z|\leq1.
\]
If $\phi(z_0)=0$ for some $|z_0|\leq1$, the absolute convergence of
$\psi(z_0)$ would imply $\theta(z_0)=0$, contradicting the assumption that
$\phi$ and $\theta$ have no common zeroes. Thus $\phi(z)\neq0$ throughout
$|z|\leq1$, and the displayed identity also gives
$\psi(z)=\theta(z)/\phi(z)$. This relation uniquely determines all of the
coefficients $\{\psi_j\}$.
\end{proof}

\begin{definition}[Invertibility]
An ARMA$(p,q)$ process defined by the equations
$\phi(B)X_t=\theta(B)Z_t$ is said to be \emph{invertible} if there exists
a sequence of constants $\{\pi_j\}_{j\geq0}$ such that
\[
\sum_{j=0}^{\infty}|\pi_j|<\infty
\]
and
\[
Z_t=\sum_{j=0}^{\infty}\pi_jX_{t-j},
\qquad t=0,\pm1,\ldots.
\]
\end{definition}

\begin{theorem}
Let $\{X_t\}$ be an ARMA$(p,q)$ process for which the polynomials
$\phi(\cdot)$ and $\theta(\cdot)$ have no common zeroes. Then $\{X_t\}$
is invertible if and only if $\theta(z)\neq0$ for all $z\in\mathbb{C}$
such that $|z|\leq1$. The coefficients $\{\pi_j\}$ in the invertible
representation $Z_t=\sum_{j=0}^{\infty}\pi_jX_{t-j}$ are determined by
the relation
\[
\pi(z)=\sum_{j=0}^{\infty}\pi_jz^j
=\frac{\phi(z)}{\theta(z)},
\qquad |z|\leq1.
\]
The coefficients $\{\pi_j\}$ can be calculated from recursion relations
analogous to those for $\{\psi_j\}$.
\end{theorem}
\begin{proof}
\end{proof}
